% ============================================================
% SCHEDE PG AGGIUNTIVE — CASA DORIA
% Personaggi per tavolo alternativo / quinta sessione
% Ruoli non coperti dai 5 PG standard
% ============================================================

\chapter{Schede PG Aggiuntive --- Casa Doria}
\index{schede PG!Doria!aggiuntivi}

% ════════════════════════════════════════════════════════════
% PG 6 — CORSO FERRANDO detto IL MASTRO
% @rpg-schema: <https://rpg-schema.org/instances/genova1507/pg-corso-ferrando>
%   a fitd:Character ;
%   rdfs:label "Corso Ferrando detto il Mastro" ;
%   fitd:crew <.../crew-doria> ;
%   fitd:role "Il Mastro d'Ascia / Il Tecnico dell'Arsenale" ;
%   fitd:action [ fitd:actionName "Tinker" ; fitd:dots 4 ] ;
%   fitd:action [ fitd:actionName "Wreck" ; fitd:dots 3 ] ;
%   fitd:action [ fitd:actionName "Survey" ; fitd:dots 2 ] .
% ════════════════════════════════════════════════════════════

\pgheader[DoriaBlue]{Corso Ferrando detto il Mastro}{Il Mastro d'Ascia / Il Tecnico dell'Arsenale}{Doria}
\pgquote{Una galea è legno, ferro e corda. So dove taglia il legno, dove cede il ferro, dove si strappa la corda. Vale per le navi. Vale per tutto il resto.}

\section*{Background \& Motivazione}

\begin{rulebox}{}
  \textbf{Background:} Quarantotto anni, capo dei calafatori e mastro d'ascia
  dell'arsenale privato dei Doria a Prè. Non è un nobile, non è un soldato —
  è l'uomo che tiene in acqua le galee. Conosce ogni nave del porto di Genova
  per nome, per età, per difetti strutturali. Sa quali chiglie marciscono sotto
  la vernice fresca e quali scafi porterebbero una tempesta atlantica.
  I francesi lo hanno cercato di corrompere tre volte. Ha detto no tre volte.
  Non per lealtà astratta — perché i Doria pagano puntuale e non fanno domande.

  \medskip\noindent\textcolor{DoriaBlue}{\textbf{Motivazione:}} L'arsenale è suo.
  Non sulla carta — nella pratica. Ci ha passato trent'anni. L'idea che un
  governatore straniero possa decidere cosa si costruisce o si demolisce in
  quello spazio lo rende fisicamente incapace di dormire.
\end{rulebox}

\section*{Attributi \& Azioni}


\begin{attributoblock}[DoriaBlue]{INSIGHT — Mente}
  \azione{Caccia}{Hunt}{1}{Trova difetti strutturali, materiali nascosti, vizi occulti}
  \azione{Studio}{Study}{2}{Tecnica navale, matematica pratica, contratti di cantiere}
  \azione{Osserva}{Survey}{2}{Valuta una nave o un edificio a colpo d'occhio}
  \azione{Ingegno}{Tinker}{4}{Costruisce, ripara, sabota, modifica qualsiasi cosa}
\end{attributoblock}


\begin{attributoblock}[DoriaBlue]{PROWESS — Corpo}
  \azione{Finezza}{Finesse}{2}{Mani precise da artigiano — anche con un coltello}
  \azione{Ombra}{Prowl}{1}{Si muove nell'arsenale e nel porto invisibile}
  \azione{Combattimento}{Skirmish}{2}{Ascia da lavoro, catena, corpo robusto}
  \azione{Forza bruta}{Wreck}{3}{Sfonda, demolisce, smonta con metodo}
\end{attributoblock}


\begin{attributoblock}[DoriaBlue]{RESOLVE — Spirito}
    \azione{Percezione}{Attune}{1}{Sente quando qualcosa non torna in un cantiere}
    \azione{Comando}{Command}{2}{Gli operai lo seguono senza discutere}
    \azione{Relazioni}{Consort}{2}{Artigiani, portuali, fornitori di materiali}
    \azione{Persuasione}{Sway}{1}{Diretto, concreto, persuasivo come un preventivo}
\end{attributoblock}

\medskip
\begin{pgclockbox}{Stress \hfill \clock{9}{0}}
  \small\textit{Traumi: $\square$ Freddo \quad $\square$ Duro \quad $\square$ Ossessionato \quad $\square$ Paranoico \quad $\square$ Irrazionale \quad $\square$ Spezzato}
\end{pgclockbox}

\section*{Abilità Speciale}

% @rpg-schema: <.../pg-corso-ferrando>
%   fitd:specialAbility [ rdfs:label "Vizi Strutturali" ] .
\begin{doriabox}{Vizi Strutturali}
  Corso conosce i difetti nascosti di ogni nave, edificio o struttura che abbia
  visto almeno una volta. Può dichiarare un \textbf{Vizio}: una debolezza specifica
  di una struttura presente nella scena. Sfruttarla non richiede tiro se si ha
  il tempo di agire con metodo; richiede \textsc{Forza Bruta} in situazioni di fretta.

  \medskip\noindent\textit{Una volta per scena. Il GM può richiedere che Corso
  abbia avuto accesso alla struttura in precedenza.}
\end{doriabox}

\section*{Legami}

\begin{table}[h]
\centering\renewcommand{\arraystretch}{1.4}
\begin{tabularx}{\linewidth}{@{}L{2.4cm} X@{}}
  \toprule\rowcolor{DoriaBlue}
  \textcolor{white}{\textbf{PG}} & \textcolor{white}{\textbf{Legame}} \\
  \midrule
  Ottaviano & L'Ammiraglio sa che senza di me le sue galee restano a secco. Lo sa e non lo dice mai. \\
  \rowcolor{DoriaBlue!8}
  Bianca & L'unica Doria che è mai scesa nell'arsenale a fare domande vere. Le ho risposto. \\
  Giacomo & Il contrabbandiere usa i miei magazzini per depositi non ufficiali. Lo so. Non dico niente. \\
  \rowcolor{DoriaBlue!8}
  Eleonora & La vedova dell'ammiraglio. Mi ha chiesto di tenerle pronta una barca. Non ho fatto domande. \\
  \bottomrule
\end{tabularx}
\end{table}

\section*{Equipaggiamento}
\begin{goldlist}
  \item Ascia da mastro — strumento di lavoro, non arma, ma difficile da distinguere
  \item Planimetria dell'arsenale con i tunnel di drenaggio (non su carta ufficiale)
  \item Lista dei materiali difettosi nelle galee francesi nel porto — compilata di notte
  \item Chiavi di tutti i magazzini dell'arsenale, inclusi tre che ufficialmente non esistono
\end{goldlist}


\begin{secretbox}
  \textbf{Segreto:} Corso ha già sabotato parzialmente una delle galee francesi nel
  porto — non abbastanza da farla affondare subito, ma abbastanza da renderla
  inaffidabile in una tempesta o in uno scontro prolungato. Lo ha fatto per
  iniziativa propria, senza ordini. Se i Doria lo scoprono, potrebbero non
  gradire l'autonomia. Se i francesi lo scoprono, è un uomo morto.

  \medskip\noindent\textcolor{SecretRed}{\textbf{Agenda:}}
  \textit{Garantire che il Colpo Riservato abbia successo anche se qualcosa va
  storto in mare. Ha preparato tre piani di contingenza tecnica. Non le ha
  dette a nessuno perché nessuno gliele ha chieste — e lui non offre informazioni
  che non vengono richieste.}
\end{secretbox}


\clearpage
% ════════════════════════════════════════════════════════════
% PG 7 — ELEONORA DORIA vedova di Simone Cattaneo
% @rpg-schema: <https://rpg-schema.org/instances/genova1507/pg-eleonora-doria>
%   a fitd:Character ;
%   rdfs:label "Eleonora Doria vedova Cattaneo" ;
%   fitd:crew <.../crew-doria> ;
%   fitd:role "La Vedova dell'Ammiraglio / Il Potere Informale" ;
%   fitd:action [ fitd:actionName "Consort" ; fitd:dots 4 ] ;
%   fitd:action [ fitd:actionName "Sway" ; fitd:dots 3 ] ;
%   fitd:action [ fitd:actionName "Survey" ; fitd:dots 3 ] .
% ════════════════════════════════════════════════════════════

\pgheader[DoriaBlue]{Eleonora Doria ved.\ Cattaneo}{La Vedova dell'Ammiraglio / Il Potere Informale}{Doria}
\pgquote{Mio marito comandava dodici galee. Io comandavo mio marito. Fate i conti.}

\section*{Background \& Motivazione}

\begin{rulebox}{}
  \textbf{Background:} Cinquantatré anni, cugina di Ottaviano, vedova del capitano
  Simone Cattaneo morto in battaglia cinque anni fa. Ha gestito sola il patrimonio
  di famiglia, i crediti commerciali, e le relazioni con i capitani che erano stati
  fedeli a suo marito. In teoria non ha autorità formale nella gerarchia dei Doria.
  In pratica, nessun capitano della flotta prende decisioni importanti senza
  consultarla. I francesi la ignorano perché è una donna. Questo è il loro errore
  più grave.

  \medskip\noindent\textcolor{DoriaBlue}{\textbf{Motivazione:}} Ha perso il marito
  in una guerra che non era sua. Non perderà anche Genova. E non lasceà che la
  rivolta venga gestita da uomini che pensano che il coraggio si misuri solo
  in battaglie navali.
\end{rulebox}

\section*{Attributi \& Azioni}



\begin{attributoblock}[DoriaBlue]{INSIGHT — Mente}
  \azione{Caccia}{Hunt}{1}{Trova persone attraverso la rete dei capitani}
  \azione{Studio}{Study}{2}{Contabilità, diritto successorio, contratti commerciali}
  \azione{Osserva}{Survey}{3}{Legge le dinamiche di una stanza con precisione chirurgica}
  \azione{Ingegno}{Tinker}{1}{Codici, lettere cifrate}
\end{attributoblock}


\begin{attributoblock}[DoriaBlue]{PROWESS — Corpo}
  \azione{Finezza}{Finesse}{2}{Gestualità calibrata, presenza fisica nei salotti}
  \azione{Ombra}{Prowl}{1}{Sa muoversi senza farsi notare in ambienti aristocratici}
  \azione{Combattimento}{Skirmish}{0}{Non combatte — non ne ha bisogno}
  \azione{Forza bruta}{Wreck}{0}{No}
\end{attributoblock}



\begin{attributoblock}[DoriaBlue]{RESOLVE — Spirito}
  
    \azione{Percezione}{Attune}{2}{Sente la lealtà e il tradimento nello stesso modo}
    \azione{Comando}{Command}{3}{Autorità materna e militare insieme — difficile da ignorare}
    \azione{Relazioni}{Consort}{4}{Capitani, armatori, mogli di nobili, preti di porto}
    \azione{Persuasione}{Sway}{3}{Convince con la memoria e il debito emotivo}
  
\end{attributoblock}

\medskip
\begin{pgclockbox}{Stress \hfill \clock{9}{0}}
  \small\textit{Traumi: $\square$ Freddo \quad $\square$ Duro \quad $\square$ Ossessionato \quad $\square$ Paranoico \quad $\square$ Irrazionale \quad $\square$ Spezzato}
\end{pgclockbox}

\section*{Abilità Speciale}

% @rpg-schema: <.../pg-eleonora-doria>
%   fitd:specialAbility [ rdfs:label "La Vedova Ricorda" ] .
\begin{doriabox}{La Vedova Ricorda}
  Eleonora ha mantenuto i rapporti con ogni capitano che ha servito sotto
  suo marito. Una volta per sessione può chiamare uno di questi capitani
  a obbedire a un ordine diretto — anche contro gli ordini formali della
  famiglia Doria, se può invocare un debito personale.

  \medskip\noindent\textit{Il capitano obbedisce senza tiro. Se l'ordine lo mette
  in pericolo fisico grave, richiede \textsc{Persuasione} da Posizione Controllata.}
\end{doriabox}

\section*{Legami}

\begin{table}[h]
\centering\renewcommand{\arraystretch}{1.4}
\begin{tabularx}{\linewidth}{@{}L{2.4cm} X@{}}
  \toprule\rowcolor{DoriaBlue}
  \textcolor{white}{\textbf{PG}} & \textcolor{white}{\textbf{Legame}} \\
  \midrule
  Ottaviano & Mio cugino. Lo rispetto più di qualunque uomo in questa città. Anche per questo non gli dico tutto. \\
  \rowcolor{DoriaBlue!8}
  Bianca & Mia nipote di cuore, non di sangue. L'ho guardata crescere. So cosa diventerà prima di lei. \\
  Fra' Matteo & Il cappellano sa cose che io so. Ci siamo riconosciuti. Non ne parliamo. \\
  \rowcolor{DoriaBlue!8}
  Corso & Il Mastro tiene pronta la mia barca. Non gli ho detto perché. Lui non ha chiesto. \\
  \bottomrule
\end{tabularx}
\end{table}

\section*{Equipaggiamento}
\begin{goldlist}
  \item Libro contabile personale con i crediti dovuti da otto capitani genovesi
  \item Anello con sigillo del marito — ancora riconosciuto come autorità da chi l'ha servito
  \item Barca a remi pronta a Sampierdarena, preparata dal Mastro
  \item Lettera di suo marito, scritta prima della battaglia. Non l'ha mai aperta. La porta sempre.
\end{goldlist}


\begin{secretbox}
  \textbf{Segreto:} Eleonora sa che Ottaviano ha ipotecato le galee. Lo ha scoperto
  attraverso i suoi canali commerciali prima ancora che Fra' Matteo lo sapesse per
  via confessionale. Ha già trovato un acquirente alternativo per il debito —
  un mercante genovese di simpatie popolari che ricomprerebb il debito a condizioni
  meno pericolose. Ma muoversi senza il consenso di Ottaviano significherebbe
  ammettere di averlo sorvegliato.

  \medskip\noindent\textcolor{SecretRed}{\textbf{Agenda:}}
  \textit{Proteggere le galee da una strada diversa da quella di Ottaviano.
  Se il Colpo Riservato riesce, il problema si risolve da solo. Se fallisce,
  Eleonora agisce autonomamente — e pagherà il prezzo politico dopo.}
\end{secretbox}

